\section{Introducción}\label{sec:introduccion}

La siniestralidad vial constituye un problema de salud pública en el Perú y en el mundo \cite{who_road_safety_2023}. Dentro de los siniestros con desenlace fatal existe una distinción operativa relevante: la mayoría deja un solo fallecido, pero una fracción (los siniestros \emph{multifatales}) concentra dos, cinco o hasta decenas de víctimas en un solo evento, típicamente asociados a transporte interprovincial, vías rurales y condiciones adversas. El Observatorio Nacional de Seguridad Vial (ONSV) del Ministerio de Transportes y Comunicaciones publica una base a nivel de registro de los siniestros fatales del país \cite{onsv_siniestros_fatales_2021_2025}, con variables de contexto como zona rural/urbana, red vial, condición climática, característica de la vía, perfil longitudinal, superficie de calzada y coordenadas geográficas.

La pregunta de investigación que guía este trabajo es: \emph{dado que ocurre un siniestro fatal, ¿qué condiciones del entorno y del evento se asocian a que sea de alta letalidad (dos o más fallecidos), y puede una red neuronal clasificarlas?}

En rigor técnico, el problema abordado es de \textbf{clasificación binaria} entre letalidad simple y multifatalidad. La capa de salida de la red produce una probabilidad continua de multifatalidad, y la categoría estimada resulta de aplicar sobre esa probabilidad un umbral de decisión validado; ambos productos se reportan por separado a lo largo del informe, porque responden a preguntas distintas: cuán probable es el desenlace y qué clase se declara. El desarrollo abarca el ciclo completo de trabajo: análisis exploratorio, preprocesamiento, ingeniería de características, entrenamiento y selección de arquitectura, calibración, evaluación, explicabilidad e interfaz gráfica de consulta.

Este encuadre condicional (\emph{dentro} de los siniestros fatales) es una decisión metodológica deliberada: la base ONSV registra únicamente siniestros con al menos un fallecido, por lo que no existe clase negativa ``no fatal''. En lugar de mezclar fuentes con criterios de registro distintos ---lo que introduciría sesgo de selección---, se reformuló el problema hacia la pregunta que esta base sí puede responder con rigor. La hipótesis subyacente es que las variables pre-impacto disponibles (mecánica del siniestro, ruralidad, clima, geometría de la vía, momento del día) contienen señal suficiente para que una red neuronal supere a una línea base trivial y compita con métodos clásicos.
